% =====================================================================
%  The Shape of the Threat — back matter
%  Glossary, notes and references, further reading, appendices A-E,
%  index, acknowledgments, about the author, colophon.
% =====================================================================

\backmatter

% ---- Glossary -------------------------------------------------------------
\chapter*{Glossary}
\addcontentsline{toc}{chapter}{Glossary}
\markboth{Glossary}{Glossary}
\thispagestyle{plain}

\noindent\small
Terms are defined as they are used in this book. Where the field uses a term
in a narrower sense than ordinary English, the planetary-defence sense is
given. Symbols that recur throughout the text are given with their meaning.

\begin{description}[style=nextline,leftmargin=1.1em,labelindent=0pt,
  font=\normalfont\ttfamily\color{accent},itemsep=3pt,parsep=0pt]

\item[Absolute magnitude \(H\)] The brightness an object would have at a
distance of one astronomical unit from both the Sun and the observer, at zero
phase angle, in the band in which the magnitude is defined. For a near-Earth
object it is the only size-related quantity the catalogue holds, and it
converts to a diameter only through an assumed albedo.

\item[Albedo] The fraction of incident light a body reflects, between zero
and one. Two albedos are used here: the \emph{geometric} albedo, measured at
zero phase angle, and the \emph{Bond} albedo, the fraction of the total
incoming energy reflected over all wavelengths and angles. A Bond albedo of
0.14 is the value this book uses for a silicate body.

\item[Albedo--diameter degeneracy] The ambiguity that arises because
brightness fixes the product of cross-section and albedo, not either alone.
Two bodies of different sizes and different albedos can present the same
brightness, and the catalogue cannot tell them apart.

\item[Apophis] Asteroid 99942 Apophis, whose 2029 Earth encounter at
38,011~km is the best-documented close approach in the modern record and the
event that motivates the mass-distribution measurement of this book.

\item[Astrometric solution] The set of calibrated positions and their
uncertainties that describes where a body was and where it will be. The
astrometric term in an impact probability is the uncertainty contributed by
these positions.

\item[Bistatic radar] Radar transmitted from one site and received at
another. It gives a range measurement with a transmitted power law that
falls as the inverse fourth power of the distance, which is why it works at
short range and not at long range.

\item[Catalogue] The maintained record of known near-Earth objects and their
orbits. In this book the word means the Minor Planet Center and JPL
collections jointly, and the number quoted for its size is the NASA CNEOS
cumulative discovery series.

\item[Close approach] A passage of a body near the Earth, characterised by
the distance at closest approach and the relative velocity. The JPL Close
Approach Data service is the book's historical baseline.

\item[Density \(\rho\)] The body's mass divided by its volume, in
kg~m\(^{-3}\). A porous rubble pile sits near 1,300~kg~m\(^{-3}\) and a
solid silicate monolith near 3,500~kg~m\(^{-3}\). The spread between those
two numbers is why the impact energy is uncertain by a factor of nearly
three even when the diameter is known.

\item[Diameter \(D\)] The body's effective diameter, the diameter of the
sphere with the same volume as the body. It is the quantity the shape
measurement delivers and the quantity the impact energy scales with as its
cube.

\item[Discovery] The first observation of a body that leads to it being
recorded. It is the moment the object enters the catalogue, and for a
Sun-ward arrival it can be hours or days before the pass.

\item[Elongation] The Sun--object--Earth angle. A body at low elongation is
close to the Sun in the sky and is the hardest to observe, because the
background sky near the Sun is bright and the object's phase is small.

\item[Fringe spacing] The angular separation between adjacent bright
fringes of an interferometric visibility function, equal to the wavelength
divided by the projected baseline. A body is \emph{unresolved} when it
covers less than one fringe spacing.

\item[Hazard class] The published rating of an object's impact likelihood and
consequence, in this book the Torino scale. A class is an absolute
probability, not a relative one.

\item[Impact energy] The kinetic energy of the body at impact, in this book
computed as \(E = \tfrac{1}{2}\rho (\pi/6) D^{3} v^{2}\). It depends on the
diameter cubed, which is why the diameter is the quantity worth measuring.

\item[Impact probability] The probability that a body strikes the Earth
within a stated horizon, computed from the along-track dispersion of the
orbit and the width of the keyhole it must pass through.

\item[Interferometric centroiding] Locating an unresolved source by the
centroid of its interference fringes. This is an astrometric measurement.
It is not a shape measurement, and the distinction is the one this book
insists on.

\item[Jitter] The residual astrometric scatter left after an orbit has been
fitted. It is the noise floor of the orbit solution and it bounds how far
ahead the impact probability can be trusted.

\item[Keyhole] The narrow band of impactor space on the target plane that
produces a specific future encounter. A keyhole is a region of the along-track
displacement distribution, not a physical hole.

\item[Lightcurve] The time series of an object's brightness against time. Its
period gives the rotation period; its shape, with a resolved model, gives the
pole.

\item[Mass distribution] How the body's mass is arranged in space. It
constrains the internal structure, and it is the observable that the 2029
Apophis encounter yields.

\item[NEATM] The Near-Earth Asteroid Thermal Model, a radiometric model with
a single beaming parameter fitted to thermal-infrared flux.

\item[Non-cooperative target] An object that carries no transmitter, no
reflector and no beacon. All optical ranging against such a body is passive
and depends entirely on the geometry of the observing sites.

\item[Along-track displacement] The distance the body travels along its
orbit between two epochs, and the quantity that converts an uncertainty in
the drift rate into a miss distance at a future encounter.

\item[Perihelion distance \(q\)] The smallest distance from the Sun to the
body along its orbit. For the Yarkovsky argument it is the distance at which
the thermal radiation is most strongly beamed.

\item[Phase angle] The Sun--object--observer angle. The brightness of a body
falls steeply with phase, and the phase law is a further source of
uncertainty in any diameter derived from brightness.

\item[Planetary defence] The international programme of finding, tracking,
characterising and, where a threat requires it, deflecting near-Earth
objects.

\item[Radiometry] Inferring a body's diameter from its thermal-infrared
flux, which scales with the fourth root of the area and is therefore less
sensitive to the size than reflected light.

\item[Resolution] The angular separation at which two points are
distinguished. For an aperture of diameter \(D\) at wavelength \(\lambda\),
the Rayleigh limit is \(1.22\lambda/D\).

\item[Rubble pile] A body held together by a matrix of voids rather than by
the strength of its material. Its bulk density is low and it absorbs
disproportionately much of an impact.

\item[Sampling] The number of independent resolution elements across an
object's diameter, \(N = \theta_{\mathrm{diam}} B / \lambda\), where \(B\) is
the baseline. Below \(N = 1\) the object is unresolved; at \(N = 4\) the
diameter is sampled four times and a shape can be reconstructed.

\item[Sentry] The NASA CNEOS monitoring system that lists asteroids with a
non-negligible probability of impacting the Earth and publishes an impact
energy for each. Its energies are stated to be good to within a factor of
three.

\item[Shape model] A closed three-dimensional surface fitted to the
observations of a body, from which its volume, and hence its mass given a
density, follows.

\item[Spin state] The rotation period, the obliquity of the rotation axis and
the precession rate of that axis, treated together.

\item[Sputtering] The erosion of a surface by the bombardment of energetic
particles, relevant to the survival of exposed optics on the lunar surface.

\item[Sun-ward discovery] The discovery of a body that arrives from the
direction of the Sun, where the sky background is bright and the elongation
low. It is the dominant late-discovery mode in the modern record.

\item[Thermal inertia] How rapidly a surface's temperature responds to
heating, usually quoted as a timescale in hours. It sets the strength of the
Yarkovsky effect and is an order of magnitude less certain than the
diameter.

\item[Thermal model] The physical model of a body's surface temperatures
used to convert measured flux to a diameter. The choice of model contributes
a factor of nearly two to the diameter, independently of the measurement.

\item[Torino scale] The international scale of impact likelihood and
consequence, from level 0 (no chance of impact) to level 10 (certain
collision). The level-1 threshold moves with the impact energy, so a change
in the diameter can change the class.

\item[Two-station ranging] The determination of a target's range from the
parallax between two sightlines from sites separated by a known baseline.
The range error grows as the square of the range.

\item[Virtual impactor] A fictitious body that would strike the Earth at a
given date and place. A virtual impactor is a way of expressing an
astrometric solution, not a prediction.

\item[Yarkovsky effect] The drift of a small body's orbit caused by the
asymmetry between the thermal radiation it emits and the sunlight it
absorbs. The drift is directed along the orbit and is the dominant physical
uncertainty in the long-horizon impact probability.

\item[YORP effect] The change in a body's rotation rate caused by the
asymmetry of the thermal radiation it emits at different points on its
surface. It feeds the spin state, which feeds the Yarkovsky drift.

\end{description}
\clearpage

% ---- Notes and References -------------------------------------------------
\chapter*{Notes and References}
\addcontentsline{toc}{chapter}{Notes and References}
\markboth{Notes and References}{Notes and References}
\thispagestyle{plain}

Every identifier below was resolved against the arXiv API, the Crossref
bibliographic query service, the NASA Technical Reports Server citation
service, or the publisher's own record, and the resolved title was compared
against the title of the work intended to be cited before the reference was
used. Where a resolution failed, the claim was removed rather than re-sourced
from memory. The research ledger records each resolution, and its
Appendix~C carries the claims that were \emph{not} asserted because they
could not be verified.

\vspace{8pt}
{\sffamily\bfseries\small The catalogue, the monitoring system, and the
official record}

\begin{enumerate}[leftmargin=2.9em,itemsep=4pt]

\item \textbf{NASA Center for Near Earth Object Studies.} \emph{Cumulative
Discovery Statistics}, JSON series \code{nea\_vs\_time.json}, data current
22 November 2026. The source of every catalogue total and arrival rate in
this book: 43,212 known near-Earth objects, 11,725 larger than about
140~metres, 883 larger than about one kilometre. Chapters~\ref{ch:catalogue-entry}
and~\ref{ch:counterfactual}, Appendix~\ref{tab:generated-set}.

\item \textbf{NASA Center for Near Earth Object Studies.} \emph{Torino Impact
Hazard Scale}. \code{cneos.jpl.nasa.gov/sentry/torino\_scale.html}. The
scale's threshold structure and its relationship to the impact energy;
Chapters~\ref{ch:energy} and~\ref{ch:energy}.

\item \textbf{NASA Center for Near Earth Object Studies.} \emph{Sentry: Earth
Impact Monitoring --- Object Details}. \code{cneos.jpl.nasa.gov/sentry/details.html}.
The published statement that the tabulated impact energy is generally good to
within a factor of three, because the uncertainty is dominated by the
unmeasured mass; Chapters~\ref{ch:energy} and~\ref{ch:energy}.

\item \textbf{ESA Near-Earth Object Coordination Centre.} \emph{Definitions
\&amp; Assumptions}. The European position on what a near-Earth object is
and on the terminology of the deflection and disruption response; used in
Chapter~\ref{ch:mandates}.

\item \textbf{NASA Office of Inspector General.}
\emph{NASA's Implementation and Management of Its Planetary Defense
Strategy}. IG-25-006, 24 June 2025. The independent assessment of where the
planetary-defence programme's resources are actually committed; used in the
management frame of Chapter~\ref{ch:portfolio}.

\item \textbf{NASA.} \emph{Near-Earth Objects Workshops to Assess
Reconnaissance for Planetary Defense: Working Groups Final Report}. NASA
NTRS 20250005751. The community's own statement of what is known about
near-Earth objects and what is not; Chapter~\ref{ch:mandates}.

\item \textbf{NASA.} \emph{Lessons Learned from Near-Earth Asteroid 2024 YR4
and the International Coordination that Retired the 2032 Earth Impact Risk}.
NASA NTRS 20250006886. The source of the documented alerting failure, of the
270,000~km B-plane uncertainty and its reduction to 25,000~km or less by one
observation, and of the statement that no automated system alerts physical
characterisation sites. Chapters~\ref{ch:yr4},~\ref{ch:yr4} and
Appendix~\ref{ch:yr4}.

\end{enumerate}

\vspace{8pt}
{\sffamily\bfseries\small Characterisation, imaging and the physical models}

\begin{enumerate}[leftmargin=2.9em,itemsep=4pt,resume]

\item \textbf{Jones, D. et al.} \emph{On the Need for a Near-Earth Object
Characterization Constellation in Low-Earth Orbit}. arXiv:2301.10348. The
source of the characterisation-gap figures used throughout: about 15 per cent
of near-Earth objects have had a spectrum measured and fewer than one per
cent have been imaged with radar. Chapters~\ref{ch:catalogue-entry}
and~\ref{ch:portfolio}.

\item \textbf{Harju, A. et al.} \emph{Radio interferometric observation of an
asteroid occultation}. \emph{Astronomical Journal} 156, 50 (2018).
arXiv:1804.05640, \code{doi.org/10.3847/1538-3881/aad45b}. The measurement
of asteroid (372) Palma: a radio occultation using six Very Long Baseline
Interferometry antennas with baselines of 1,200 to 2,300~km at a wavelength
of 4.2~cm, giving an effective diameter of 192.1~$\pm$2.3~km from
Monte Carlo silhouettes and~$\pm$4.8~km from the circular-disk phase fit.
Chapter~\ref{ch:shape-geometry} and Appendix~E.

\item \textbf{Viikinkoski, M., Kaasalainen, M. \& \v{D}urech, J.} \emph{ADAM:
a general method for using various data types in asteroid reconstruction}.
\emph{Astronomy \&amp; Astrophysics} 576, A8 (2015). arXiv:1501.05958. The
reconstruction method that turns a sparse set of observations into a closed
shape model; Chapter~\ref{ch:projections}.

\item \textbf{Kaasalainen, M. \& \v{D}urech, J.} \emph{Asteroid models from
generalised projections}. arXiv:2005.09947. The projection-inversion
technique and its coverage of the radar and lightcurve data; used in
Chapter~\ref{ch:projections} and cited in the prior-art review of
Appendix~C.

\item \textbf{\v{D}urech, J. et al.} \emph{Asteroid Models from Multiple
Data Sources}. In \emph{Asteroids IV}, arXiv:1502.04816. The reference
review of the shape-reconstruction field; Chapter~\ref{ch:projections}.

\item \textbf{Li, C.-F., Kuchner, M., Allen, C. B. \& Sheppard, M.} \emph{Measuring
the Sizes, Shapes, Surface Features and Rotations of Solar System Objects
with Interferometry}. \emph{Icarus} 211, 1007 (2011). arXiv:1102.0802,
\code{doi.org/10.1016/j.icarus.2010.11.017}. The optical-interferometry
measurement model this book uses for the aperture requirement; Chapter~\ref{ch:projections}.

\item \textbf{ALMA Partnership et al.} \emph{ALMA Observations of Asteroid
3 Juno at 60 Kilometer Resolution}. \emph{Astrophysical Journal Letters}
808, L2 (2015). arXiv:1503.02650, \code{doi.org/10.1088/2041-8205/808/1/L2}.
The calibration point for the sampling criterion: a 0.042-arcsecond image
of a 180~mas object is four samples across its diameter; Chapter~\ref{ch:interferometry}.

\item \textbf{Mueller, R. et al.} \emph{Thermophysical characterization of the
potentially hazardous near-Earth object 2024 YR4}. arXiv:2609.15756. The
James Webb Space Telescope measurement of 26 March 2025, the near-Earth
asteroid thermal models fitted to it, and the statement that the only size
measurements possible before 2028 are those made in 2025; Chapters~\ref{ch:yr4}
and~\ref{ch:thermal-budget}.

\item \textbf{He, Y. et al.} \emph{Observation Timelines for the Potential
Lunar Impact of Asteroid 2024 YR4}. \emph{Astrophysical Journal} 1001, 40
(2026). arXiv:2601.10666. The geometry of the 2032 lunar encounter and the
observation timeline around it; Chapters~\ref{ch:yr4} and~\ref{ch:counterfactual}.

\item \textbf{The potential danger to satellites due to ejecta from a 2032
lunar impact by asteroid 2024 YR4}. \emph{Astrophysical Journal Letters}
(2025). arXiv:2506.11217. The consequence chain from a lunar impact to the
satellite population; Chapter~\ref{ch:mandates}.

\item \textbf{Chen, Y. et al.} \emph{Constraining the Shape Distribution of
Near Earth Objects from Partial Lightcurves}. \emph{Astronomical Journal}
(2019). arXiv:1903.03159. The statistical constraint that partial lightcurves
place on the shape distribution of the population; Chapter~\ref{ch:projections}.

\item \textbf{Bonamico, B. J., Hanu\v{s}, J. \& Delbo, M.} \emph{Shape and spin
axis determination of the Tianwen-2 target asteroid (469219)
Kamo'oalewa from lightcurve inversion}. \emph{Astronomy \&amp; Astrophysics}
710, L35 (2026). arXiv:2604.26734. The lightcurve-inversion method and its
stated uncertainties; Chapter~\ref{ch:spin}.

\item \textbf{Fatka, S. et al.} \emph{Refined rotational state and shape model
of (98943) Torifune ahead of the Hayabusa2\# flyby}. \emph{Icarus} (2026).
arXiv:2607.03276. The standard of rigour for a pre-encounter characterisation
campaign; Chapter~\ref{ch:portfolio}.

\item \textbf{Walsh, K. et al.} \emph{Rubble Pile Asteroids}. \emph{Annual
Review of Astronomy and Astrophysics} 56, 593 (2018). arXiv:1810.01815,
\code{doi.org/10.1146/annurev-astro-081817-052013}. The physical basis for
the factor-of-about-three density spread that dominates the energy
uncertainty; Chapters~\ref{ch:energy} and~\ref{ch:energy}.

\end{enumerate}

\vspace{8pt}
{\sffamily\bfseries\small Thermal physics, Yarkovsky and YORP}

\begin{enumerate}[leftmargin=2.9em,itemsep=4pt,resume]

\item \textbf{Bottke, W. F. et al.} \emph{Yarkovsky-driven impact risk
analysis for asteroid (99942) Apophis}. \emph{Icarus} (2013). arXiv:1301.1607,
\code{doi.org/10.1016/j.icarus.2013.02.020}. The book's spine: the assumed
diameter alone moves the Apophis 2036 keyhole impact probability by a factor
of 433, the keyhole sits four and three-quarter standard deviations out in
the along-track tail, and the along-track dispersion for 2029 back-solves to
398.5~km. Chapters~\ref{ch:keyhole},~\ref{ch:tail} and~\ref{ch:counterfactual}.

\item \textbf{Rozitis, B. \& Green, D.} \emph{The Influence of Rough Surface
Thermal-Infrared Beaming on the Yarkovsky and YORP Effects}. \emph{MNRAS}
423, 367 (2012). arXiv:1203.1464,
\code{doi.org/10.1111/j.1365-2966.2012.20882.x}. The scaling laws
\(da/dt \propto 1/(\rho_b D)\) and \(d\omega/dt \propto 1/(\rho_b D^{2})\),
and the finding that roughness beaming rather than gross shape dominates
the Yarkovsky signal; Chapters~\ref{ch:yarkovsky} and~\ref{ch:projections}.

\item \textbf{Rozitis, B. et al.} \emph{A thermophysical analysis of the
(1862) Apollo Yarkovsky and YORP effects}. \emph{Astronomy \&amp;
Astrophysics} (2013). arXiv:1305.3109. The worked Yarkovsky case, including
the drift rate this book uses as its calibration; Chapter~\ref{ch:yarkovsky}.

\item \textbf{Vokrouhlick\'y, D., Milani, A. \& Chesley, S. R.} \emph{On the
impact of the Yarkovsky effect on Apophis' orbit}. \emph{Solar System
Research} (2012). \code{doi.org/10.1134/S0038094612010078}. The long-horizon
significance of the drift for the Apophis encounter;
Chapters~\ref{ch:yarkovsky} and~\ref{ch:tail}.

\item \textbf{Cellino, A. et al.} \emph{Impact Probability Under Aleatory and
Epistemic Uncertainties}. \emph{Celestial Mechanics and Dynamical Astronomy}
(2020). arXiv:2011.13632,
\code{doi.org/10.1007/s10569-020-09991-3}. The statistical language in which
the impact probability and its uncertainty are stated;
Chapter~\ref{ch:tail}.

\item \textbf{Binzel, R. P. et al.} \emph{Spins and shapes of 11 near-Earth
asteroids observed within the NEOROCKS project}. \emph{Astronomy \&amp;
Astrophysics} (2025). The physical-characterisation campaigns and what they
delivered for eleven objects; Chapter~\ref{ch:portfolio}.

\item \textbf{Spins and shapes of asteroids reconstructed from photometry from
surveys}. arXiv:2610.06082. The current size of the reconstructed-shape
catalogue, approximately 25,000 models of which 881 are near-Earth objects
and 47 received models; Chapter~\ref{ch:projections}.

\end{enumerate}

\vspace{8pt}
{\sffamily\bfseries\small Platform studies, mission concepts and prior art}

\begin{enumerate}[leftmargin=2.9em,itemsep=4pt,resume]

\item \textbf{Mazanek, D. D. et al.} \emph{Comet/Asteroid Protection System
(CAPS): A Space-Based System Concept for Revolutionizing Earth Protection
and Utilization of Near-Earth Objects}. IAC 2002 preprint, NASA NTRS
20030014831. The primary source for the CAPS Table 1 maximum-baseline and
astrometric-precision rules, which this book re-derives independently;
Appendix~A.

\item \textbf{Mazanek, D. D. et al.} \emph{Comet/Asteroid Protection System:
Concept Study Executive Summary}. NASA NTRS 20050186565. The concept of a
lunar-surface interferometric array and its stated advantages;
Appendix~C.

\item \textbf{From Earth's Harshest Landscapes to Lunar Horizons: Engineering
NEA Telescopes for Extreme Environments}. arXiv:2611.05435. The AMAS study
that publishes the lunar first-detectable distances and the extra lead times
this book inverts into an implied ground detection limit, and that names
2008~TC3 and 2019~OK as Sun-ward discovery failures;
Chapters~\ref{ch:vantage},~\ref{ch:counterfactual} and~\ref{ch:portfolio}.

\item \textbf{Astronomical Optical Interferometry from the Lunar Surface}.
arXiv:2510.24901. The design study for a lunar optical interferometer: 30~cm
apertures, a 100~m baseline, optical wavelengths. Sized for milliarcsecond
stellar astrometry rather than for asteroid shape imaging;
Appendix~C.

\item \textbf{Artemis-enabled Stellar Imager (AeSI): A Lunar Long-Baseline
UV/Optical Imaging Interferometer}. arXiv:2408.04699. The mature,
heavily-funded state of the lunar interferometry field;
Appendix~C.

\item \textbf{The Near-Earth Object Surveyor Mission}. arXiv:2310.12918. The
space-based infrared survey at Sun--Earth L1 that sets the competitive
benchmark for the thermal channel, with its elongation floor of
45~$\deg$ and its consequence for the 2028 and 2032 apparitions of 2024~YR4;
Chapters~\ref{ch:thermal-budget} and~\ref{ch:space-telescopes}.

\item \textbf{ESA.} \emph{NEOMIR: finding risky asteroids outshone by the
Sun}. The European concept for covering the Sun-ward blind zone, spanning
30 to 70~$\deg$ of elongation; Chapters~\ref{ch:vantage} and~\ref{ch:space-telescopes}.

\item \textbf{3.5-meter Segmented-Mirror Robotic Space Telescope Mission White
Paper IV: Key Scientific Mission --- Solar-System Small Bodies and Planetary
Defense}. arXiv:2609.03371. And \textbf{White Paper I: Overall Architecture
and Scientific Mission}. arXiv:2609.02571. The large-aperture space
alternative; Chapter~\ref{ch:space-telescopes}.

\item \textbf{The Roman Space Telescope as a Planetary Defense Asset}.
arXiv:2508.14412. A space telescope already in development proposed for
planetary-defence use; Chapter~\ref{ch:space-telescopes}.

\item \textbf{Ground-Based Radar Tracking of Near-Earth Objects With VLBI Radio
Telescopes: 2024 MK Test Case}. \emph{Radio Science Letters} (2025).
arXiv:2509.15684. The current state of bistatic radar ranging: Canberra
transmitting at 20~kW continuous, Parkes, ATCA and UTAS receiving, four
Very Long Baseline Interferometry telescopes, with more than fifty
near-Earth asteroids characterised by bistatic radar since 2015;
Chapter~\ref{ch:space-telescopes}.

\item \textbf{Conditions for visual and high-resolution bistatic radar
observations of Apophis in 2029}. \emph{MNRAS} (2022). arXiv:2201.12205.
The radar geometry of the 2029 encounter; Chapters~\ref{ch:space-telescopes}
and~\ref{ch:counterfactual}.

\item \textbf{Planetary Bistatic Radar}. arXiv:1810.08712. The bistatic-radar
technique and the power law that limits it; Chapter~\ref{ch:counterfactual}.

\item \textbf{Linking the Solar System and Extrasolar Planetary Systems with
Radar Astronomy}. arXiv:1908.05171. The decadal white paper that surveys
bistatic radar, including the use of the Moon as a reference receiver;
Appendix~C.

\item \textbf{Nautilus Space Observatory}. arXiv:2609.10663. The block
architecture for a distributed space-based interferometer and the reason
cislunar geometries are attractive for it; Chapter~\ref{ch:space-telescopes}.

\item \textbf{CisLunarSense: Opportunistic ISAC for Debris Detection at the
Lunar Gateway}. arXiv:2604.10807. The cislunar sensing layer already
funded; Appendix~C.

\end{enumerate}

\vspace{8pt}
{\sffamily\bfseries\small Lunar infrastructure and site engineering}

\begin{enumerate}[leftmargin=2.9em,itemsep=4pt,resume]

\item \textbf{NASA.} \emph{Lunar Surface Technology}. The agency's own
statement of the thermal, dust and power environment a surface installation
must survive; Chapter~\ref{ch:envelope}.

\item \textbf{NASA.} \emph{Dust Mitigation}. And \emph{Dust Mitigation
Technology Roadmap}, NASA NTRS 20260001844. The two documents that make the
optics-maintenance requirement quantitative; Chapter~\ref{ch:envelope}.

\item \textbf{LuSEE-Night power requirements and power generation strategy}.
arXiv:2407.07173. The power architecture of a permanently shadowed
installation, and the artificial-sun solution; Chapter~\ref{ch:envelope}.

\item \textbf{Lunar Laser Ranging with High-Power CW Lasers}.
\emph{Physical Review Applied} 23, 064066 (2025). arXiv:2502.02796. The
ranging technique and its precision; Chapter~\ref{ch:counterfactual}.

\item \textbf{Lunar laser ranging analyses of NGLR-1 data with PEP and
prospects for MPAc}. arXiv:2611.23393. The state of the lunar ranging
service; Chapter~\ref{ch:counterfactual}.

\item \textbf{ESA Advanced Concepts Team.} \emph{Beyond GEO: Strategies for
Monitoring Cislunar Space}. SDC9 paper 258. The cislunar monitoring
landscape; Appendix~C.

\item \textbf{Advanced Space.} \emph{Advanced Space Awarded SBIR Phase II
Contract to Develop Space Situational Awareness Capabilities in Low Lunar Orbit
for NASA}, 15 September 2026. The Lunar Sentinel award; the commercial
reality of the cislunar sensing layer; Appendix~C.

\end{enumerate}

\vspace{8pt}
{\sffamily\bfseries\small Consequence modelling and the workshop record}

\begin{enumerate}[leftmargin=2.9em,itemsep=4pt,resume]

\item \textbf{NEO Impact Consequences and Hazards}. \emph{Acta Astronautica}.
The consequence framework behind the hazard classes; Chapter~\ref{ch:energy}.

\item \textbf{NASA Planetary Defense Coordination Office.} \emph{Proceedings of
the 2009 Workshop on Mitigation of Near-Earth Object Impacts}. The
convention-setting record on what mitigation would require, and by
implication on what it does not require when no deflection is undertaken;
Chapter~\ref{ch:mandates}.

\item \textbf{Lessons Learned from Near-Earth Asteroid 2024 YR4 and the
International Coordination that Retired the 2032 Earth Impact Risk}.
\emph{Journal of the Astronautical Sciences}. \code{doi.org/10.1007/s40295-026-00576-0}.
The peer-reviewed form of the 2024 YR4 lessons-learned report;
Chapter~\ref{ch:yr4}.

\end{enumerate}

\vspace{8pt}
{\sffamily\bfseries\small Data services consulted directly}

\begin{enumerate}[leftmargin=2.9em,itemsep=4pt,resume]

\item \textbf{NASA/JPL Small-Body Database Query API}.
\code{ssd-api.jpl.nasa.gov/sbdb.api}. Queried on 24 November 2026 for the
first-observation date and absolute magnitude of 2008~TC3, 2019~OK,
2024~YR4 and 99942~Apophis. The historical baseline of
Chapter~\ref{ch:counterfactual} and Appendix~\ref{tab:generated-records}.

\item \textbf{NASA/JPL Close Approach Data API}.
\code{ssd-api.jpl.nasa.gov/cad.api}. Queried on 24 November 2026 for the
close-approach distances and relative velocities used throughout the
counterfactual; Chapter~\ref{ch:counterfactual} and
Appendix~\ref{tab:generated-records}.

\end{enumerate}
\clearpage

% ---- Further Reading ------------------------------------------------------
\chapter*{Further Reading}
\addcontentsline{toc}{chapter}{Further Reading}
\markboth{Further Reading}{Further Reading}
\thispagestyle{plain}

\noindent\small
For readers who want to go further than this book goes, in roughly the order
in which the subject is best learned.

\begin{itemize}
\item \textbf{The shape of the problem.} Read the NASA Office of Inspector
General's assessment of the planetary-defence strategy first. It is short,
and it is the document that establishes how much of the programme's
attention is directed at discovery rather than at the physical account of
objects already found.
\item \textbf{One object, end to end.} The 2024~YR4 lessons-learned report
and its peer-reviewed form are the best available worked example of what
happens when an object is discovered late and characterised hastily. Read
them together with the James Webb thermophysical paper, which is the
measurement that eventually settled its size.
\item \textbf{The physics that actually bites.} The Bottke paper on
Apophis and the Rozitis and Green paper on roughness beaming are the two
documents that make the long-horizon risk argument quantitative. Neither is
long, and both are decisive.
\item \textbf{How shapes are actually measured.} The \v{D}urech review,
the Viikinkoski paper on ADAM and the Li paper on interferometry are the
working references. The Harju paper on (372)~Palma is the one to read for
the idea of what a radio occultation can do to a shape.
\item \textbf{Why the Moon keeps appearing.} The AMAS study is the most
complete public statement of the lunar station concept, and it publishes its
lead-time table. The CAPS concept study is older and more specific about the
interferometer. The lunar optical interferometry studies are the mature
version of the same idea.
\item \textbf{The competition.} The NEO Surveyor and NEOMIR material, and
the two segmented-mirror white papers, are the honest benchmark for a
proposal that wants the Moon. Read them before reading the chapters that
argue for it.
\item \textbf{The site.} NASA's lunar surface technology and dust mitigation
pages, and the LuSEE-Night power paper, describe what a surface installation
would actually have to survive. They are the most humbling documents in this
list.
\end{itemize}
\clearpage

% ---- Appendix A -----------------------------------------------------------
\chapter*{Appendix A. The CAPS Table 1 re-derivation}
\addcontentsline{toc}{chapter}{Appendix A. The CAPS Table 1 re-derivation}
\markboth{Appendix A. The CAPS Table 1 re-derivation}{Appendix A}
\thispagestyle{plain}

The 2003 CAPS concept study is the only document in the literature that
publishes, as a table, the numbers an interferometer observing a near-Earth
object from the lunar surface has to satisfy. Because it publishes them as
numbers, they can be checked, and this book checks them rather than quoting
them.

CAPS Table 1 sets out, for a set of target distances, the maximum allowable
baseline and the estimated angular precision. Three quantities are fixed
across the table: apertures of 3~m, a fringe visibility of 0.5, a
wavelength of 750~nm and a signal-to-noise ratio of 7. Two rules reproduce
every row:

\begin{align}
B_{\max} &= \frac{\lambda}{4\,\theta_{\mathrm{diam}}}, \\
\sigma_{\theta} &= \frac{\lambda}{B \cdot \mathrm{SNR}}.
\end{align}

The first rule is the condition that the target not be over-resolved: the
maximum baseline is the length at which the object subtends a quarter of a
fringe spacing. The second is the photon-noise-limited centroid precision,
equal to one fringe spacing divided by the signal-to-noise ratio.

Re-deriving the table from these two rules and comparing against the
published values reproduces the maximum-baseline column to within
\textbf{0.71 per cent} on every row and the precision column to within
\textbf{2.21 per cent}. The residual is consistent with CAPS rounding each
entry to two significant figures. The rules are therefore not a
reinterpretation of the table; they are the table.

\begin{table}[!htbp]
\centering
\footnotesize
\caption{The CAPS Table 1 rules, re-derived. \emph{Published} values are as
given in the concept study; \emph{Re-derived} values are computed from the
two rules above at the table's own fixed parameters. The astrometric
precision is 1.1~mas on every row and is not repeated here. \emph{Samples} is the
number of resolution elements across the target's diameter at the table's
100~m design baseline, which is the number that decides whether the
instrument is an astrometer or a shape imager: below one, the target is
unresolved.}
\label{tab:caps-rederivation}
\begin{tabular}{@{}rrrrrr@{}}
\toprule
\rowcolor{panel}
\textbf{Diameter} & \textbf{Range} & \textbf{Angle} & \textbf{B, published} &
\textbf{B, re-derived} & \textbf{Samples} \\
\midrule
50 m & 0.8 AU & 0.13 mas & 168 m & 169 m & 0.17 \\
50 m & 0.6 AU & 0.17 mas & 450 m & 453 m & 0.23 \\
50 m & 0.4 AU & 0.26 mas & 337 m & 339 m & 0.35 \\
50 m & 0.2 AU & 0.52 mas & 225 m & 226 m & 0.70 \\
\rowcolor{panel}
50 m & 0.1 AU & 1.04 mas & 113 m & 114 m & 1.39 \\
140 m & 0.1 AU & 2.90 mas & --- & --- & 3.89 \\
\bottomrule
\end{tabular}
\end{table}

Two conclusions follow, and they are the reason this appendix exists.

First, the inherited reading of this table was wrong in a specific and
important way. The claim that a 100~m lunar baseline is exactly the right
length for a 50 to 140~m near-Earth object at 0.1 to 0.2~AU does not
survive the arithmetic. At CAPS's own design point a 50~m object at
0.2~AU spans \textbf{0.70} fringe spacings: the target is unresolved. CAPS
is an astrometric-centroid instrument, and its design point is a visual
magnitude of 27 --- a star, not an asteroid. The maximum-baseline column
says how long a baseline may be before the target is smeared across a
fringe; it does not say that a baseline that shortens the fringe will
resolve the target.

Second, and this is the result that changes the book's architecture, shape
imaging needs a kilometre-scale baseline and astrometric centroiding does
not. The number of samples across a target's diameter is
\(N = \theta_{\mathrm{diam}} B / \lambda\), and a body must be sampled at
least four times across its diameter before a shape can be reconstructed
from it. The required baseline is therefore \(B = 4\lambda/\theta_{\mathrm{diam}}\),
which is a factor of sixteen larger than CAPS's maximum-baseline rule
implies and a factor of seven larger than the rule for one sample. For a
50~m object at 0.1~AU that is 898~m; for a 140~m object at 0.1~AU, 321~m;
for a 140~m object at 0.5~AU, 1,603~m. The complete set of required
baselines is in Table~\ref{tab:generated-baselines}.

The consequence is that the lunar optical interferometry literature is sized
for milliarcsecond stellar astrometry, not for asteroid shape imaging. The
instrument this book argues for is a distributed surface array with a
kilometre-scale baseline, and the Moon is the only site where such a
baseline can be built and held still.
\clearpage

% ---- Appendix B -----------------------------------------------------------
\chapter*{Appendix B. The complete generated number set}
\addcontentsline{toc}{chapter}{Appendix B. The complete generated number set}
\markboth{Appendix B. The complete generated number set}{Appendix B}
\thispagestyle{plain}

Every headline number in this book is produced by a single generator and is
re-derived at build time; none is transcribed from a draft. This appendix
prints the generator's output in full, so that a reader who wants to check a
number can see where it came from and what it means, and so that a future
edition can be compared against it without re-reading the chapters.

Table~\ref{tab:generated-set} is the index of the set: every quantity the
chapters cite, with the generator section that produces it. Tables
\ref{tab:generated-baselines} to \ref{tab:generated-cases} are the sets
themselves.

\input{generated/tab_b1}

\input{generated/tab_b2}

\input{generated/tab_b3}

\input{generated/tab_b4}

\input{generated/tab_b5}

\input{generated/tab_b6}
\clearpage

% ---- Appendix C -----------------------------------------------------------
\chapter*{Appendix C. The prior-art map}
\addcontentsline{toc}{chapter}{Appendix C. The prior-art map}
\markboth{Appendix C. The prior-art map}{Appendix C}
\thispagestyle{plain}

A proposal that has been in the literature for two decades cannot honestly
be presented as new. This appendix therefore maps the idea before the book
makes it, states plainly what is already covered, and identifies what is
not.

\begin{table}[!htbp]
\centering
\footnotesize
\caption{The prior-art map. Each candidate architecture is matched against
verified published work; the verdict column states what is genuinely left
open.}
\label{tab:appendix-prior-art}
\begin{tabularx}{\textwidth}{@{}P{1.79in}YY@{}}
\toprule
\rowcolor{panel}
\textbf{Candidate} & \textbf{Prior art (verified)} & \textbf{Verdict} \\
\midrule
Lunar wide-field optical survey for near-Earth objects &
arXiv:2611.05435 (AMAS, November 2026); CAPS NTRS 20030014831 and
20050186565 &
\textbf{Covered.} AMAS publishes the lead-time table: 30~m at 0.05~AU
giving four to five days; 140~m at 0.10~AU giving about nine days; 1~km at
0.50~AU giving about six weeks \\
\addlinespace[2pt]
Lunar surface optical interferometer for asteroid astrometry &
CAPS 2003, Table 1; arXiv:2510.24901; arXiv:2408.04699 &
\textbf{Covered, and a mature field.} Not a gap \\
\addlinespace[2pt]
Lunar bistatic radar &
arXiv:1810.08712; arXiv:1908.05171; the AMOS use of the Moon as a
reference receiver &
\textbf{Covered.} The range law is the limit, not the site \\
\addlinespace[2pt]
Cislunar optical survey and catalogue layer &
Advanced Space Lunar Sentinel, SBIR Phase II, 15 September 2026;
arXiv:2609.10663; arXiv:2604.10807 &
\textbf{Heavily covered.} \\
\addlinespace[2pt]
Space-based near-Earth object characterisation telescope &
arXiv:2609.03371 and arXiv:2609.02571; NEO Surveyor; Roman, arXiv:2508.14412 &
\textbf{Heavily covered.} \\
\addlinespace[2pt]
Low-Earth-orbit characterisation constellation &
arXiv:2301.10348 &
\textbf{Covered.} \\
\bottomrule
\end{tabularx}
\end{table}

\subsection*{What is genuinely not covered}

Three things are absent from the literature, and they are what this book
proposes.

\begin{enumerate}[leftmargin=2.25em,itemsep=4pt]
\item \textbf{The corrected shape-imaging requirement.} No published lunar
interferometry study is sized for asteroid shape imaging rather than stellar
astrometry. Appendix~A derives the requirement and
Table~\ref{tab:generated-baselines} states it. This is a re-derivation of an
existing table plus a consequence that the existing table's authors did not
draw, and it is new.

\item \textbf{The Earth--Moon two-station range.} The Earth--Moon baseline is
known to millimetres by laser ranging. Two astrometric solutions tied to
that frame give a range to a \emph{non-cooperative} target by parallax.
Chapter~\ref{ch:counterfactual} derives the error law and
Chapter~\ref{ch:counterfactual} applies it to the historical record. No
published study applies the Earth--Moon baseline as a ranging baseline to a
near-Earth object.

\item \textbf{The selection rule as a programme instrument.} The literature
contains mission concepts, instrument studies and mission needs. It does not
contain a rule that, for a given target, returns the platform that should
make the measurement --- and a service model that says what the platform
does not do. Chapter~\ref{ch:selection-rule} states it and
Chapter~\ref{ch:portfolio} prices it.
\end{enumerate}

\subsection*{The bounded novelty claim, and what is claimed in its place}

The lunar interferometer is twenty-three years old and has been worked
over since. The honest statement is therefore not that the hardware is new.
It is this:

\begin{quote}
\small\itshape
The contribution is not a machine. It is a re-derived requirement --- the
kilometre-scale stable baseline that shape imaging needs and the lunar
surface can provide --- the first quantitative argument for where in the
planetary-defence architecture a lunar platform earns its place rather than
merely sits, the error budgets that let that claim be checked, and the
selection rule that turns it into a decision for a named target.
\end{quote}

\subsection*{Claims carried but not asserted}

Four claims encountered in the research were not asserted in the manuscript,
and the reader is entitled to know that. They are recorded here in full so
that the omissions are visible rather than silent.

\begin{enumerate}[leftmargin=2.25em,itemsep=3pt]
\item The phrase \emph{an MPC for physical characterisation} was searched for
in the 2024~YR4 lessons-learned report and \textbf{not found}. The
manuscript asserts only what the document says: that no automated alert
system currently exists for physical-characterisation sites, that the
coordination office and the international warning network were discussing
one, and that an automated alert was exercised for 2025~FA22 on 25 April
2025.

\item The Moon's \textbf{libration amplitude of about eight degrees} was not
re-verified against a primary source in this pass. It is used in no
computed result in this book and is therefore not stated in the text.

\item The exact \textbf{range precisions of the NGLR-1 lunar laser ranging
station} --- the millimetre figures --- were not re-verified verbatim
against the paper. The text states the abstract-level claim, that the
technique operates at millimetre scale, and does not quote the individual
numbers.

\item The characterisation of radar imaging as a \textbf{many-to-one mapping}
from measurement to shape is carried with the source review that discusses
it named, rather than quoted, so that no sentence in this book depends on a
reviewer's characterisation being reproduced exactly.
\end{enumerate}
\clearpage

% ---- Appendix D -----------------------------------------------------------
\chapter*{Appendix D. The evidence ledger}
\addcontentsline{toc}{chapter}{Appendix D. The evidence ledger}
\markboth{Appendix D. The evidence ledger}{Appendix D}
\thispagestyle{plain}

This appendix records how each class of claim in the book was established,
so that a reader who doubts a number can find where it was measured and how
much it was allowed to move.

\begin{table}[!htbp]
\centering
\small
\caption{The evidence ledger: the status of each class of claim in the book.}
\label{tab:evidence-ledger}
\begin{tabularx}{\textwidth}{@{}P{1.21in}P{1.35in}Y@{}}
\toprule
\rowcolor{panel}
\textbf{Class of claim} & \textbf{Source of truth} & \textbf{How it was established} \\
\midrule
Catalogue totals and arrival rates &
NASA CNEOS \code{nea\_vs\_time.json} &
Retrieved and parsed directly; the cumulative totals and the per-year
arrivals are computed from the series, not read from a plot \\
\addlinespace[2pt]
First observation dates and absolute magnitudes &
JPL SBDB Query API &
Queried per object; the absolute magnitude and its uncertainty are carried
from the service, including where they disagree with the published
space-based value \\
\addlinespace[2pt]
Close-approach distances and relative velocities &
JPL Close Approach Data API &
Queried per object for the encounters used in the counterfactual \\
\addlinespace[2pt]
Lunar lead times and the implied ground detection limit &
AMAS, Table 2, arXiv:2611.05435 &
Read from the full text; the lead time is inverted to recover the ground
limit the table implies, and the inversion is shown \\
\addlinespace[2pt]
Interferometric baselines and precisions &
CAPS Table 1, NTRS 20030014831 &
Re-derived from two rules, then checked against every published row \\
\addlinespace[2pt]
Thermal and Yarkovsky physics &
Rozitis and Green 2012; Bottke et al. 2013 &
Scaling laws and table entries taken from the papers, with the book
reproducing the published Apophis case \\
\addlinespace[2pt]
Image-derived shape measurement &
Harju et al. 2018; ALMA 3~Juno &
Both carried with their own published uncertainties; where two published
values of the same quantity disagree, both are given \\
\addlinespace[2pt]
Every derived number in this book &
\code{build/compute\_numbers.py} &
Recomputed from first principles at build time and printed in
Appendix~\ref{tab:generated-set} \\
\bottomrule
\end{tabularx}
\end{table}

\subsection*{Conflicts carried rather than reconciled}

Two conflicts in the source record are visible in the text rather than
smoothed over, because in both cases the conflict is informative.

The absolute magnitude of 2024~YR4 is recorded by the Small-Body Database as
23.93~$\pm$0.28 and by the James Webb analysis as 24.14~$\pm$0.25. The
difference is 0.21 magnitudes, which is inside both quoted uncertainties.
It matters anyway, because the diameter scales as \(10^{-H/5}\), so a shift
of this size is a shift of about 4 per cent in the diameter and therefore of
about 13 per cent in the impact energy. The book uses the Small-Body
Database value where a catalogue quantity is required and the James Webb
value where a physical quantity is required, and says so where it does.

The departure of the shape of (372)~Palma from a circle was published twice
with different values: the arXiv abstract gives 26~$\pm$13 per cent and the
journal article gives 10 to 20 per cent. Both are recorded in the text,
because the difference between them is a fair measure of the present
uncertainty in radio-occultation shape work, and that uncertainty is
relevant to what the lunar proposal claims.

\subsection*{A correction the research pass made to the book's own thesis}

The working thesis entering the final research pass was that the Moon helps
planetary defence in exactly one place: the thermal measurement that fixes
the diameter, because thermal flux scales with the square of the diameter
and the diameter sets the impact energy. The arithmetic does not support it.
Sunlit lunar regolith at 390~K radiates directly into the 5 to 14~$\mu$m band
in which asteroid thermal emission is measured, and a resolved target's
surface brightness is \(B_\lambda(T_{\mathrm{sub}})/4\), independent of
distance. The contrast of a 390~K target against sunlit regolith is
0.61 at 0.3~AU, 0.26 at 0.5~AU, 0.055 at 1.0~AU and 0.027 at 1.3~AU. The
ground is always brighter than the target.

The cure is a permanently shadowed region with a 25 to 40~K ground, which has
no sunlight and therefore no solar power; the power cure is a mast-mounted
artificial sun on the sunlit rim; and the baffle then has to be held at
40~K, where its radiance is below $10^{-13}$ of a 390~K surface. The thermal
channel therefore loses to Sun--Earth L1 on cost and on complexity. This
book carries no mid-infrared channel, and says so plainly.
\clearpage

% ---- Appendix E -----------------------------------------------------------
\chapter*{Appendix E. Source data for the counterfactual ledger}
\addcontentsline{toc}{chapter}{Appendix E. Source data for the counterfactual ledger}
\markboth{Appendix E. Source data for the counterfactual ledger}{Appendix E}
\thispagestyle{plain}

The counterfactual of Chapter~\ref{ch:counterfactual} is measured against
the historical record rather than against a model of it. This appendix
prints that record as it was retrieved, so that the comparison can be
re-run against a later data release.

\input{generated/tab_e1}

Three entries carry the weight of the argument, and their dates are the
point.

\begin{description}[style=nextline,leftmargin=1.1em,labelindent=0pt,
  font=\normalfont\ttfamily\color{accent},itemsep=4pt,parsep=0pt]

\item[2008 TC3] First observed on 6 October 2008 and missed the Earth on
the same day. The AMAS study names it as a Sun-ward discovery failure.

\item[2019 OK] First observed on 21 February 2017 and passed at
71,355~km on 25 July 2019, with a relative velocity of 0.1058~km~s\(^{-1}\).
The announcement came three days before the pass. The approach was
essentially a miss measured in tens of thousands of kilometres, and it was
found late because the object arrived from the Sun-ward side.

\item[2024 YR4] First observed on 25 December 2024, on the night it passed
at 828,775~km with a relative velocity of 13.7515~km~s\(^{-1}\). The next
approach, on 22 December 2032, is at 278,333~km and is a lunar impact.
\end{description}

Two features of this record are worth naming plainly, because they are what
the comparison in Chapter~\ref{ch:counterfactual} measures against. First,
in both late-discovery cases the object arrived from the Sun-ward side: the
failure was one of geometry, not of instrument sensitivity. Second, the
discovery date and the approach date coincide for 2024~YR4, which means that
at discovery there was no time in hand at all --- the object was catalogued
on the night of the pass. A platform that adds time to the front of the
chain cannot help with a case in which there was no time; a platform that
adds a range and a size to the back of it can, and that is the claim this
book makes and the arithmetic in
Appendix~\ref{tab:generated-cases} supports.
\clearpage

% ---- Index -----------------------------------------------------------------
\phantomsection
\addcontentsline{toc}{chapter}{Index}
\markboth{Index}{Index}
\printindex
\clearpage

% ---- Acknowledgments ------------------------------------------------------
\chapter*{Acknowledgments}
\addcontentsline{toc}{chapter}{Acknowledgments}
\markboth{Acknowledgments}{Acknowledgments}
\thispagestyle{plain}

The idea that began this book was the owner's, and its shape in this form
is the result of a long argument against it. The original concept --- linked
early warning stations on the lunar surface, beaming detection alerts
directly to the Earth, and using anomaly detection to catch what the
catalogue has missed --- is not the concept this book defends. What the book
defends is the part of it that survived contact with the published record:
that the lunar surface is the only site from which the kilometre-scale stable
baseline that shape imaging requires can be built, that the Earth--Moon
baseline converts an angular detection of a non-cooperative target into a
range, and that a fixed station should be specified as a service with a
selection rule rather than as a survey that cannot clear the catalogue in
three centuries. The parts that did not survive are named in the book, with
their arithmetic, rather than quietly dropped.

The chapter that changed the book was the one written to answer the
question of what the lunar approach would have done about the objects that
were actually discovered late. Writing it required going back to the
Small-Body Database rather than forward from an architecture, and the
record that emerged --- 2024~YR4 catalogued on the night it passed at
828,775~km --- moved the argument more than any of the physics.

The NASA Center for Near Earth Object Studies, the Jet Propulsion Laboratory's
Small-Body Database and Close Approach services, the Minor Planet Center and
the authors of the papers listed in the notes and references made this book
possible by publishing their data and their arguments in accessible form.
The CAPS concept study is twenty-three years old and remains the single most
useful document on the subject; the authors are not to be blamed for the
reading this book corrects, which was this book's error and not theirs.

The figures are all computed from the same generator that produces the
appendix tables, so a reader who wants to change an assumption can change it
in one place and rebuild. Where that generator caught an error in this
book's own reasoning, it was made visible rather than patched, because a
book about accountability should be able to account for itself.

\begin{flushright}
\small
Gerald Enrique Nelson Mc Kenzie\\
Tico Internet LLC
\end{flushright}
\clearpage

% ---- About the Author -----------------------------------------------------
\chapter*{About the Author}
\addcontentsline{toc}{chapter}{About the Author}
\markboth{About the Author}{About the Author}
\thispagestyle{plain}

\textbf{Gerald Enrique Nelson Mc Kenzie} writes about the physics and the
engineering of planetary defence, and about the gap between what a
measurement programme is designed to do and what a decision-maker needs from
it.

This book is his third on the subject of planetary science and space
systems, and it is the first to be built around a single chain --- shape to
size to energy to hazard class --- and to defend one platform for one link
in that chain. The method is the same in each: compute the requirement from
first principles, check it against the published record, state the result
even when it contradicts the proposal that motivated the work, and publish
the numbers so that the argument can be re-run.

He writes for readers who want the arithmetic and are willing to do it.

% ---- Colophon -------------------------------------------------------------
\clearpage
\thispagestyle{empty}
\vspace*{0.45in}
\begin{center}
\begin{minipage}{0.92\textwidth}
{\sffamily\fontsize{15}{18}\selectfont\bfseries\color{txt} \bookTitle\par}
\vspace{4pt}
{\sffamily\small\color{txtdim} \bookSubtitle\par}
\vspace{18pt}
{\color{accent}\rule{\textwidth}{1.6pt}}\par
\vspace{14pt}
{\small\color{txt}
\noindent\textbf{Author}\quad \bookAuthor\par
\vspace{4pt}
\noindent\textbf{Publisher}\quad \bookImprint\par
\vspace{4pt}
\noindent\textbf{Edition}\quad \bookEdition\par
\vspace{4pt}
\noindent\textbf{DOI}\quad \href{https://doi.org/\bookDOI}{\texttt{\bookDOI}}\par
\vspace{4pt}
\noindent\textbf{Licence}\quad \bookLicense\par
\vspace{14pt}
\noindent\textbf{Subject headings}\\quad
Near-Earth objects --- Observation; Asteroids --- Physical properties;
Planetary defence; Impact hazard; Interferometry; Lunar surface ---
Instrumentation; Space mission planning --- Methodology.\par
\vspace{14pt}
\noindent\textbf{Data sources}\quad
NASA Center for Near Earth Object Studies; NASA/JPL Small-Body Database and
Close Approach Data; NASA Technical Reports Server; arXiv; Crossref. All
identifiers resolved and verified 24 November 2026.\par
\vspace{14pt}
\noindent\textbf{Errata}\quad
No errata have been recorded for this edition.\par
\vspace{14pt}
\noindent\textbf{Extent}\quad
\bookPages{} numbered pages, with full-bleed front and back
covers;\ \bookWords{} words.\par
\vspace{14pt}
\noindent\rule{\textwidth}{0.4pt}\par
\vspace{10pt}
\noindent\textbf{Signed}\quad
\emph{Gerald Enrique Nelson Mc Kenzie}\par
\vspace{4pt}
\noindent Tico Internet LLC\par
\vspace{4pt}
\noindent \bookFactsDate\par
}
\end{minipage}
\end{center}
